\input _template

\setlanguage{SK}

\Title{Cifrovanie}

\Author{Patrik Bak}

\sec Úvod

Ukážeme si základné tipy a~triky používané v~úlohách s~ciframi.

\sec Techniky

\begitems \style i
\i Číslo $\overline{a_1a_2\dots a_n}$ rozpíšeme pomocou dekadického zápisu, napr. $\overline{ab}=10a+b$ alebo $\overline{abc}=100a+10b+c$.
\i Číslo so sekvenciou cifier $x$ nasledovanou sekvenciou cifier $y$ dĺžky $k$ zapíšeme ako $x \cdot 10^k + y$, napr. 4267 je pre $k=2$ rovné $42 \cdot 10^2 + 67$.
\i Pri úlohách s~ciframi sa oplatí pozerať na zvyšky po delení~3 alebo~9, lebo číslo dáva rovnaký zvyšok po delení~3 alebo~9 práve keď je tomu tak pre ciferný súčet.
\i Ďalšie dobré kritériá pre cifry sú pre deliteľnosť $2^n$ (posledných $n$ cifier tvorí číslo deliteľné~$2^n$).
\i A ešte jedno kritérium: číslo je deliteľné~11 práve keď súčet cifier na párnych miestach dáva rovnaký zvyšok po delení~11 ako súčet cifier na nepárnych miestach.
\enditems 

\sec Úlohy

\Problem{0}{}{
	Nájdite všetky štvorciferné čísla $\overline{abcd}$ s~ciferným súčtom~12 také, že $\overline{ab}-\overline{cd}=1$.
}{	
	Platí $\overline{ab}=\overline{cd}+1$. Rozoberte prípady podľa toho, či pri sčítaní napravo nastáva prechod cez~10.
}{	
}

\Problem{0}{}{
	Určte všetky prirodzené čísla $n$, pre ktoré platí
	$$
	n+p(n)=70,
	$$
	pričom $p(n)$ označuje súčin všetkých cifier čísla $n$.
}{	
	Zapíšte $n=\overline{ab}$ a~rovnicu prepíšte.
}{	
}

\Problem{0}{}{
	Dvojciferné číslo $\overline{ab}$ nazveme \textit{nafúknuteľné}, ak z neho po pripočítaní 990-násobku vhodného jednociferného čísla získame štvorciferné číslo tvaru $\overline{axxb}$ s~nenulovou cifrou~$x$. Koľko nafúknuteľných čísel existuje?
}{	
	Zapíšte definíciu nafúknuteľnosti pomocou rovnice $\overline{ab}+990y=\overline{axxb}$, ktorú ďalej upravte.
}{
	Mali by sme prísť do rovnice $9(y-a)=x$. Lenže~$x$ je nenulová cifra, to nám veľa možností nedá.
}{	
}

\Problem{0}{}{
	Zistite, aké najmenšie kladné celé číslo možno vložiť medzi dvojčíslia $20$ a $16$ tak, aby výsledné číslo bolo násobkom čísla $2016$.
}{
	Pozrime sa na modulo~$9$.
}{
	Keďže~$2016$ je deliteľné~9, tak ciferný súčet vkladaného čísla musí byť deliteľný~9, lebo ciferný súčet $20$ a~$16$ je~$9$. 
}{
	Postupne skúšajte vkladať $9, 18, \dots$.
}

\Problem{0}{}{
	Nájdite všetky trojciferné čísla, ktoré sú súčtom faktoriálov svojich cifier.
}{
	Takéto číslo nemôže mať príliš veľké cifry. Akú najväčšiu môže mať?
}{
	Rozmyslíme, že najväčšia cifra v~čísle môže byť 5. Musí tam vôbec nejaká byť?
}{
	Rozmyslíme, že cifra~5 tam nutne musí byť. Už to nie je veľa možností, obzvlášť keď sa zamyslíme nad tým, čo môže byť prvá cifra.
}{
}

\Problem{0}{}{
	Nájdite všetky štvorciferné čísla také, že sú druhou mocninou celého čísla, a~zároveň ich prvá cifra je rovnaká ako druhá a~tretia je rovnaká ako štvrtá.
}{
	Zapíšme si naše číslo ako $\overline{aabb}$ a~rozpíšme, čomu je rovné.
}{
	Vychádza nám $\overline{aabb}=11(100a+b)$. Aby bolo toto štvorec, potrebujeme, aby $100a+b$ bolo deliteľné~11.
}{
	Platí $100a+b=99a+(a+b)$, takže nutne $11 \mid a+b$. To už nie je veľa možností.
}{
}

\Problem{0}{}{
	Nájdite najväčšie prvočíslo zložené z~10 rôznych cifier.
}{
	Pozrite sa na ciferný súčet nášho čísla.
}{
}

\Problem{0}{}{
	Zdôvodnite, že každý palindróm s~párnym počtom cifier je deliteľný~11.
}{
	Kritérium deliteľnosti~11.
}{
}

\Problem{0}{}{
	Číslo $123456789101112\dots 100$ postupne nahradzujeme ciferným súčtom, až kým nezískame jednociferné číslo. Aké číslo nám zostane?
}{
	Stačí nájsť ciferný súčet a použiť deliteľnosť~9.
}{
	Na spočítanie ciferného súčtu nášho čísla sa zamyslime, koľkokrát sa každá z~cifier vyskytuje v~našom čísle.
}{
}

\Problem{0}{}{
	Existuje číslo deliteľné~11 zložené z~cifier 1, 2, 3, 4, 5 a~6 tak, že každú použijeme práve raz?
}{
	Použijeme kritérium deliteľnosti~11 -- chceme, aby rozdiel súčtu cifier na párnych miestach~$s_p$ a~na nepárnych miestach~$s_n$ bol deliteľný~11. Skúste nájsť možnosti pre tieto súčty.
}{
	Uvedomíme si, že $s_p+s_n$ je ciferný súčet nášho čísla, takže $1+2+3+4+5+6=21$. Z~toho nájdeme možnosti pre dvojice $s_p$ a~$s_n$ a~tie zanalyzujeme.
}{
}

\Problem{0}{}{
	Mysli si trojciferné číslo, ktoré neobsahuje nulu a má prvú a poslednú cifru rôznu. Napíš ho odzadu a odčítaj menšie číslo od väčšieho. Vzniknuté číslo prípadne doplň zľava nulou na trojciferné, znova prepíš odzadu a tieto dve čísla sčítaj. Aké najmenšie a~aké najväčšie číslo sme mohli dostať?
}{
	Zapíšme si číslo ako $\overline{abc}$. Čo nám zostane po prvej operácii?
}{
	V~ďalšom kroku máme $|\overline{abc}-\overline{cba}|=99|a-c|$. Vypíšte si tieto možné čísla. 
}{
}

\Problem{0}{}{
	Nájdite všetky čísla s~prvou číslicou~6 také, že po jej odstránení zostane 25-krát menšie číslo.
}{
	Zapíšte si číslo ako $6 \cdot 10^k + x$, kde $k$ je počet cifier~$x$.
}{
	Zostavte si rovnicu zo zadania: $6 \cdot 10^k + x = 25x$. Z~toho vyjadrite $x$ a~zistite, pre aké $k$ bude $x$ celé číslo.
}{
}

\Problem{0}{}{
	Dokážte, že číslo $\underbrace{44\dots4}_{n+1}\underbrace{88\dots8}_{n}9$ je pre každé celé nezáporné číslo~$n$ štvorcom prirodzeného čísla.
}{
	Rozpíšme si číslo pomocou desatinného zápisu.
}{
	Pripomeňte si, že číslo zapísané ako $k$ rovnakých cifier $c$ sa dá vyjadriť ako $c \cdot \frac{10^k - 1}{9}$. Nezabudnite každú časť čísla vynásobiť príslušnou mocninou desiatky podľa toho, koľko cifier za ňou nasleduje.
}{	
}

\Problem{0}{}{
	Je dané prirodzené číslo $n$. Určte poslednú číslicu čísla $n^2$, ak viete, že~7 je jeho predposledná číslica.
}{
	Zapíšte číslo ako $10x+y$, kde $y$ je posledná cifra. Ako vyzerá $n^2$?
}{
	Platí $n^2=100x^2+20xy+y^2$. Kedy môže byť 7 predposledná cifra?
}{
	Analyzujte vplyv jednotlivých členov súčtu $100x^2+20xy+y^2$ na miesto desiatok.
}{
	Hlavná idea je, že $100x^2$ neprispieva na miesto desiatok vôbec a~$20xy$ prispieva párnou cifrou, takže $y^2$ musí mať nepárnu cifru na mieste desiatok.
}{	
}

\Problem{1}{}{
	Existujú dve rôzne mocniny dvojky s rovnakým počtom cifier, z ktorých by sa jedna dala získať iba preusporiadaním cifier tej druhej.
}{
	Pozrime sa na modulo~$9$.
}{
	Skúmajme mocniny dvojky modulo~9. Jednotlivé mocniny $2^1, 2^2, \dots$ sú $2,4,8,7,5,1$ a~tak~ďalej. Za akých podmienok $2^a$ a~$2^b$ dávajú rovnaký zvyšok po delení~9?
}{
}

\Problem{1}{}{
	Pre ktoré prirodzené $n$ existuje $n$-ciferné číslo deliteľné $5^n$, ktoré má všetky cifry nepárne?
}{
	Skúsme také čísla zostrojovať induktívne. 
}{
    Vezmime si vyhovujúce číslo, ktoré má $n$ cifier. Dokážte, že pripojením jednej z~cifier $\{1,3,5,7,9\}$ na začiatok čísla dostaneme vyhovujúce číslo s $n+1$ ciframi.
}{
}

\Problem{1}{}{
	Uvažujme nasledovný proces: Začneme s~ľubovoľným prirodzeným číslom $a_1$ a~vygenerujeme postupnosť $a_1,a_2,\dots$ tak, že $a_{n+1}$ dostaneme z~$a_n$ prilepením cifry rôznej od~9. Dokážte, že v~tejto postupnosti bude nevyhnutne existovať nekonečne veľa zložených čísel.
}{
	Postupujeme sporom. Čo ak by tam bolo len konečne veľa zložených čísel? Rozmyslite si, aké cifry vôbec dokážeme pripájať.
}{
	Určite nevieme pripájať párne čísla a~5. Zostáva 1, 3, 7. Skúste vylúčiť niektoré z~nich.
}{
	Idea je vylúčiť 1 a~7. Tieto môžeme mať len konečne veľa kvôli tomu, že by sme poškodili deliteľnosť~3. Zostávajú trojky. Môžeme po čase mať iba tie?
}{
	Trik na trojky je rozpísať si číslo ako 
	$$
	a\cdot 10^n + \underbrace{\overline{33\dots3}}_n=
	a\cdot 10^n + \frac{10^n-1}{3}.	
	$$
	Druhú časť vieme tiež napísať explicitne. Rozmyslite si, že by stačilo, aby $10^n-1$ bolo deliteľné~$3a$.
}{
}

\Problem{2}{}{
	Dokážte, že existuje nekonečne veľa kladných celých čísel $n$ takých, že $n^2$ zapísané v~štvorkovej sústave obsahuje iba číslice $1$ a~$2$. 
}{
	Skúsme také čísla zostrojovať induktívne. 
}{
	Trikom je, že z~troch kópií čísla vieme zostrojiť nové, ktoré vyhovuje. Nebojte sa zosilniť indukčný predpoklad.
}{
}

\DisplayExerciseSolutions

\DisplayHints

\DisplayProblemSolutions

\bye